% ═════════════════════════════════════════════════════════════════════════════
\section{Regular Model Checking}
\label{sec:rmc}
% ═════════════════════════════════════════════════════════════════════════════

Regular Model Checking (RMC) is a symbolic framework for verifying
parameterized and infinite-state systems by representing both sets of
configurations and transition relations using the theory of regular languages
and finite automata~\cite{KW01,BW98,BJNT00}. The key idea is that regular sets
serve as an efficient symbolic representation of infinite state spaces,
playing a role analogous to that of Binary Decision Diagrams (BDDs) in symbolic
model checking of finite-state systems~\cite{AJNS04,Abd12,Abd21}.

\subsection{The Framework}

In its basic form, RMC represents a system as follows. A \emph{configuration}
is a finite word over a fixed alphabet~$\Sigma$, where each symbol encodes the
local state of one component (a process, a stack element, a channel message,
etc.). Sets of configurations, including the initial set and the set of bad
configurations, are represented as regular languages over~$\Sigma$. The
transition relation is encoded as a \emph{length-preserving finite-state
transducer}~$T$ over pairs $(a, a') \in \Sigma \times \Sigma$: the transducer
reads two words of the same length in parallel, letter by letter, and accepts
the pair $(w, w')$ if and only if the system can move from configuration~$w$
to configuration~$w'$ in one step.

This simple setup already captures a wide variety of models. Parameterized
systems consisting of an arbitrary number of finite-state processes in a
linear or ring topology can be modeled by letting each position in the word
represent one process. Systems operating on unbounded queues, stacks, or
integer variables can similarly be encoded by letting word positions represent
data structure contents. The length-preserving constraint on the transducer is
not a restriction for safety verification: for parameterized systems it simply
means that the number of processes is fixed at the start of a run, while
being arbitrary across runs.

\paragraph{A token-passing example.}
Consider a linear array of processes in which a single token is passed from
left to right. Each process is either holding the token or not holding it. We
write $\Tok$ and~$\NoTok$ for these two local states. The initial configurations are those in which
the leftmost process holds the token, that is, the regular language
$\Tok\NoTok^*$. In one step, the process holding the token passes it to its
right neighbor. Figure~\ref{fig:transducer} shows a transducer~$T$ for this
transition relation. It copies a prefix of processes without the token, then
reads the pair $(\Tok,\NoTok)$ for the process that gives the token away,
followed by the pair $(\NoTok,\Tok)$ for the neighbor that receives it, and
finally copies the remaining processes. The safety property states that
exactly one process holds the token, so the bad configurations form the
regular language of words with no~$\Tok$ or with at least two occurrences
of~$\Tok$.

\begin{figure}[t]
\centering
\begin{tikzpicture}[->,>={Stealth[round]},shorten >=1pt,auto,
    node distance=2.3cm,semithick,
    every state/.style={minimum size=7mm,inner sep=1pt}]
  \node[state,initial,initial text={}] (q0) {$q_0$};
  \node[state] (q1) [right=of q0] {$q_1$};
  \node[state,accepting] (q2) [right=of q1] {$q_2$};
  \path (q0) edge[loop above] node {$(\NoTok,\NoTok)$} (q0)
        (q0) edge node {$(\Tok,\NoTok)$} (q1)
        (q1) edge node {$(\NoTok,\Tok)$} (q2)
        (q2) edge[loop above] node {$(\NoTok,\NoTok)$} (q2);
\end{tikzpicture}
\caption{A length-preserving transducer for passing the token one position to
the right. Each transition reads a letter of the current configuration
together with a letter of the next configuration.}
\label{fig:transducer}
\end{figure}

\paragraph{Reachability.}
Let $\Init$ and $\Bad$ be the regular sets of initial and bad
configurations. We write $T^*$ for the reflexive transitive closure of the
transducer relation, and $\Init \circ T^*$ for the set of configurations
reachable from~$\Init$. Safety verification reduces to checking whether
$(\Init \circ T^*) \cap \Bad = \emptyset$. Liveness verification reduces to a
\emph{repeated reachability} problem: checking whether there exists an
infinite computation from~$\Init$ through~$T$ that visits a set of accepting
configurations~$F$ infinitely often.

In the token-passing example, the relation~$T^k$, obtained by composing~$T$
with itself $k$~times, moves the token exactly~$k$ positions to the right, and
each $T^k$ is recognized by a transducer with $k+2$ states. Iterating~$T$ one
step at a time therefore produces an infinite sequence of growing transducers
and never reaches a fixpoint. The transitive closure~$T^+$, however, is
regular: it moves the token any positive number of positions to the right,
and it is recognized by the transducer of Figure~\ref{fig:transducer} with an
additional loop labeled $(\NoTok,\NoTok)$ on the state~$q_1$. The reachable
set is then $\NoTok^*\Tok\NoTok^*$, which is disjoint from the bad
configurations.

\subsection{Computing Transitive Closures}

The central computational challenge is computing $T^+$ (or $\Init \circ
T^*$) effectively. Since the problem is undecidable in general, no method can
succeed on all inputs, and naively iterating the transducer, computing $T^1,
T^2, T^3, \ldots$, does not converge, as the example shows. The standard
approach is to use \emph{acceleration} or \emph{widening} techniques that jump
directly to a representation of~$T^+$.

\paragraph{Acceleration by quotienting.}
Acceleration is typically based on \emph{quotienting}: one finds an
equivalence relation~$\simeq$ on the states of a finite approximation~$T^n$,
merges equivalent states to obtain a quotient transducer $T^n/{\simeq}$, and
checks whether the result characterizes~$T^+$. The difficulty is that
quotienting generally enlarges the accepted language: $L(T^n) \subseteq
L(T^n/{\simeq})$, often strictly. The quotient is therefore taken on the
\emph{history transducer}, a canonical infinite transducer built as the
union of all $T^n$ for $n \geq 1$, where each state records the full sequence
of $T$-states visited during product construction. A key theoretical
contribution is the identification of equivalence relations on the states of
the history transducer whose quotient accepts exactly~$T^+$. Such relations
can be derived from forward and backward bisimulations, and in later work from
simulations, on these states. When the resulting quotient is finite, it is a
transducer for~$T^+$~\cite{AJNO02,AJNO03}. Related constructions iterate
transducers by identifying equivalent states with bisimulation-like
relations~\cite{DLS01}.

\paragraph{Extrapolation and abstraction.}
A different idea is to compute a few iterates $T^1, \ldots, T^k$, detect a
regular pattern of growth between consecutive iterates, and
\emph{extrapolate} it to a guess of the limit, which is then checked for
correctness~\cite{BLW03,BT02}. In the token-passing example, the transducer
for $T^{k+1}$ has one more state, and one more $(\NoTok,\NoTok)$ step between
the pair that gives the token away and the pair that receives it, than the
transducer for~$T^k$. Extrapolation turns this repeated step into the loop
on~$q_1$. \emph{Abstract regular model
checking}~\cite{BHV04} gives up exactness instead: it computes reachable sets
by alternately applying the transducer and abstracting the resulting automata,
for example by merging states that accept the same words up to a given length.
The abstraction guarantees termination, and spurious counterexamples are used
to refine it.

\paragraph{Learning regular invariants.}
For safety verification, it is not necessary to compute the reachable set
exactly: it suffices to find a regular set that contains the initial
configurations, is closed under the transducer, and is disjoint from the bad
configurations. Such a \emph{regular inductive invariant} can be searched for
with techniques from automata learning, in which a learner proposes candidate
automata and a teacher checks the three conditions and returns
counterexamples~\cite{HV05,NJ13,CHLR17}. Checking whether a given automaton is
an inductive invariant is a decidable problem on automata, so the learning
approach replaces the computation of a fixpoint by a sequence of decidable
checks.

\paragraph{Extensions.}
Beyond the linear word model, the framework extends to configurations
represented as trees, enabling verification of systems with hierarchical or
recursive structure~\cite{AJMO02,BT02}. Liveness properties can be handled by
extending the framework to infinite words~\cite{BLW05}, and regular proofs of
liveness have been studied for randomized parameterized systems~\cite{LR16}.
RMC has been successfully applied to token-passing systems, mutual exclusion
algorithms, and systems with unbounded communication channels, stacks, and
integer variables~\cite{AJNS04,Abd12}.

\subsection{Monotonic Abstraction}

While the transducer-based approach is general, it can be
computationally expensive and may not scale to larger protocols. An important
line of work within the RMC framework addresses this by replacing the full
power of regular languages with a lighter symbolic representation based on
\emph{upward-closed sets}, in a technique known as \emph{monotonic
abstraction}~\cite{ADHR07,Abd12}. The approach targets parameterized systems
with a linear process topology and global transitions (both universal and
existential). Configurations are again words over the process state alphabet,
and the subword ordering~$\preceq$ of Section~\ref{sec:wqo} plays a central
role. Bad configurations are characterized by upward-closed sets with respect
to this ordering: if a configuration is bad, then any larger configuration
(with more processes) is also bad. Such sets have a finite symbolic
representation in terms of their minimal elements, making them far more
tractable than arbitrary regular sets.

The core idea is to replace the concrete transition relation~$\to$ with an
abstract relation~$\to_A$ that over-approximates~$\to$ and is monotone with
respect to~$\preceq$. Local and existential transitions are already monotone
and require no approximation. Universal transitions, where a process may
move only if all processes in its context satisfy some condition, are
\emph{not} monotone, since adding more processes to a configuration may
violate the universal condition and disable the transition. The abstraction
handles this by allowing a configuration to silently drop processes that
obstruct the transition, effectively making the universal condition vacuously
satisfiable. Since $\to_A$ is an over-approximation, any safety property
that holds in the abstract model is guaranteed to hold in the concrete one.

Safety verification then proceeds by \emph{backward reachability analysis}, as
in Section~\ref{sec:wqo}. Starting from the upward-closed set $\Bad$ of bad
configurations, one computes the sequence $U_0 = \Bad$, $U_{i+1} = U_i \cup
\Pre_A(U_i)$, where $\Pre_A$ is the predecessor operator under~$\to_A$.
Because $\to_A$ is monotone, each $U_i$ remains upward-closed and can be
represented by a finite set of generators. Termination is guaranteed by
Higman's lemma: the subword ordering~$\preceq$ is a well-quasi-ordering on
words, so every ascending chain of upward-closed sets stabilizes after finitely
many steps. The algorithm terminates when no new generators (up to
subsumption) are produced; if none of the generators represents an initial
configuration, the system is safe.

\paragraph{The running example.}
In the running example, the abstract relation allows a process in~$\Wait$ to
move to~$\Crit$ in any configuration, after deleting the processes that are
in~$\Crit$. For instance, $\Wait\Crit \to_A \Crit$. Backward reachability
starts from the single generator $\Crit\Crit$ and computes the
configurations from which one abstract step leads to a configuration with at
least two processes in~$\Crit$. The steps from $\Idle$ to $\Wait$ and from
$\Crit$ to $\Idle$ do not create processes in~$\Crit$, so a configuration
that reaches $\upc{(\Crit\Crit)}$ by one of these steps already has at least
two processes in~$\Crit$. The abstract step from $\Wait$ to
$\Crit$ deletes all other processes in~$\Crit$, so it always results in a
configuration with exactly one process in~$\Crit$. Hence $\Pre_A$ adds no new
generators, the fixpoint is $\upc{(\Crit\Crit)}$, and since no initial
configuration $\Idle^n$ contains $\Crit\Crit$, the protocol is safe. In this
example the abstraction loses no precision with respect to the property. For
protocols in which a process inspects the others one at a time, such as
Burns' and Szymanski's algorithms, the same approach applies to the
non-atomic versions of the global conditions~\cite{ABDR08}.

\paragraph{Extensions and refinement.}
Monotonic abstraction has been extended to processes with numerical variables
whose values are compared using gap-order constraints~\cite{ADR09}, to tree
topologies, broadcast communication, and binary rendez-vous, and to systems
running under weak memory models (Section~\ref{sec:wmm}). Since the
abstraction may introduce spurious counterexamples, it has been combined with
a counterexample-guided refinement loop that restricts the deletions allowed
in the abstract transitions~\cite{ACDHHR10}. The connection between monotonic
abstraction and well-structured transition systems provides general
decidability results for the abstract systems, for classes of systems that
extend well beyond the purely parameterized setting~\cite{Abd10}.
